\chapter{Git 不是用来"保存代码"，是用来"证明历史"}
\label{ch:git_proof}

\section{故事引入：审稿意见要求复现，你却找不回当时的代码}

论文提交三个月后，审稿意见回来了。其中一条写得很直接："请提供代码和数据，我们想要复现表3的结果。"

你心里一紧——赶紧打开代码仓库。但眼前的景象让你后背发凉：

\begin{itemize}
  \item Git 历史只有寥寥几次提交："initial commit"、"update"、"fix bug"、"final version"；
  \item 论文里的结果是三个月前跑的，你已经记不清用的是哪个版本的代码；
  \item 代码目录里有多个版本的实验脚本：\texttt{train.py}、\texttt{train\_v2.py}、\texttt{train\_final.py}，你不确定用的是哪个；
  \item 更糟的是，你发现最近为了新实验已经大改了模型代码，现在的版本跑不出论文里的数字。
\end{itemize}

你只能硬着头皮回复："我们正在整理代码，会尽快提供。"然后开始痛苦的"考古工作"——试图从记忆、聊天记录、实验笔记里拼凑出当时的代码状态。

\textbf{这个场景你熟悉吗？}

\section{为什么"随手 commit"救不了你}

很多人以为自己在用 Git，但实际上只是把 Git 当成了"云盘"：

\begin{itemize}
  \item 改了一堆代码，最后一起 commit，message 随便写个"update"；
  \item 从不用分支，所有改动都在 main 上叠加；
  \item 实验跑完了才想起来 commit，这时候代码已经又改了；
  \item commit 历史里找不到"哪个版本对应哪个实验结果"。
\end{itemize}

这种用法的问题在于：\textbf{你失去了 Git 最核心的价值——作为"历史证明工具"的能力。}

在工程开发中，Git 的主要作用是协作和回滚。但在研究中，Git 的核心价值是\textbf{证明}：

\begin{itemize}
  \item 证明这个结果是用哪个版本的代码跑出来的；
  \item 证明论文里的每个实验都有对应的代码版本；
  \item 证明你可以回到任意历史版本，重新跑出相同的结果。
\end{itemize}

\section{研究中的 Git 使用陷阱}

\subsection{陷阱 1：commit 粒度太大，找不到关键改动}

\textbf{症状：}一个 commit 包含了十几个文件的改动，涉及数据处理、模型结构、训练流程等多个方面。commit message 只写了"improve model"。

\textbf{后果：}
\begin{itemize}
  \item 无法定位某个指标变化是由哪个改动引起的；
  \item 想要回滚某个错误改动，却发现无法单独撤销；
  \item 几个月后再看，完全想不起来这个 commit 做了什么。
\end{itemize}

\textbf{正确做法：}
\begin{itemize}
  \item 每个 commit 只包含\textbf{一个逻辑改动}；
  \item commit message 写清楚"改了什么"和"为什么改"；
  \item 遵循"原子性原则"：每个 commit 都应该让代码保持可运行状态。
\end{itemize}

\subsection{陷阱 2：实验和代码改动时间错位}

\textbf{症状：}先改代码跑实验，结果不错，过了两天才 commit；或者 commit 之后又临时改了几个参数重跑。

\textbf{后果：}
\begin{itemize}
  \item 产生了结果的代码版本（commit）其实对不上；
  \item 别人（包括未来的你）拿着 commit hash 复现，结果对不上；
  \item 审稿人要求复现时，你根本找不到准确的代码版本。
\end{itemize}

\textbf{正确做法：}
\begin{itemize}
  \item \textbf{先 commit，再跑实验}；
  \item 每次实验的 run.json 记录当时的 commit hash 和 dirty 状态；
  \item 如果临时改了代码，要么重新 commit，要么在 run 记录里注明 dirty 修改内容。
\end{itemize}

\subsection{陷阱 3：分支使用不当，主线混乱}

\textbf{症状：}所有实验都在 main 分支上做，各种尝试性改动和稳定代码混在一起；或者创建了很多分支但从不清理，分支之间关系混乱。

\textbf{后果：}
\begin{itemize}
  \item main 分支不稳定，充斥着各种试验性代码；
  \item 想要找到"论文复现版本"时，不知道该用哪个分支；
  \item 分支太多导致团队成员不知道该基于哪个分支开展新工作。
\end{itemize}

\section{适合研究的 Git 分支策略}

与工程项目不同，研究项目的分支策略需要平衡两个需求：
\begin{itemize}
  \item \textbf{稳定性：}论文结果必须有一个干净、稳定的代码版本支撑；
  \item \textbf{探索性：}新想法需要快速试错，不能被流程束缚。
\end{itemize}

\subsection{推荐的分支结构}

\begin{verbatim}
main (或 stable)：
  - 只接受验证通过的改动
  - 每个 merge 必须经过 DoD 检查（见第5章）
  - 确保任何时候都可以复现论文结果

exp/<hypothesis-name>：
  - 每个实验假设一条分支
  - 命名要清晰：exp/attention-ablation, exp/data-augmentation
  - 短命分支：验证完就合并或删除
  - 允许"脏"的快速迭代

archive/<paper-version>：
  - 论文提交、发表等重要节点的存档分支
  - 从 main 分支创建，不再合并回去
  - 永久保留，确保可追溯
\end{verbatim}

\subsection{典型工作流程}

\paragraph{场景 1：验证新假设}

\begin{enumerate}
  \item 从 main 创建新分支：\texttt{git checkout -b exp/new-loss-function}
  \item 在分支上快速迭代、试错，commit 可以随意
  \item 跑出有希望的结果后，整理代码
  \item 创建规范的实验记录（config + run.json）
  \item 合并回 main：\texttt{git checkout main \&\& git merge exp/new-loss-function}
  \item 删除实验分支：\texttt{git branch -d exp/new-loss-function}
\end{enumerate}

\paragraph{场景 2：论文提交}

\begin{enumerate}
  \item 确保 main 分支所有论文实验都可复现
  \item 创建存档分支：\texttt{git checkout -b archive/icml2026-v1}
  \item 在 main 上打 tag：\texttt{git tag -a paper-icml2026-v1 -m "ICML 2026 submission version"}
  \item 推送 tag：\texttt{git push origin paper-icml2026-v1}
\end{enumerate}

\paragraph{场景 3：并行探索多个方向}

\begin{enumerate}
  \item 同时创建多个实验分支：
  \begin{itemize}
    \item \texttt{exp/architecture-search}
    \item \texttt{exp/data-augmentation}
    \item \texttt{exp/loss-function}
  \end{itemize}
  \item 各分支独立推进，互不干扰
  \item 每个分支的实验产物用独立的 run\_id 管理
  \item 有价值的改动逐个合并回 main
  \item 无价值的分支直接删除
\end{enumerate}

\section{用 Tag 标记里程碑：让论文结果永久可追溯}

Tag 是 Git 中被严重低估的功能。对于研究项目，tag 的价值在于：
\begin{itemize}
  \item 为论文的每个关键版本打上永久标记；
  \item 即使 main 分支继续演进，也能精确回到历史版本；
  \item 提供清晰的版本命名，便于引用和复现。
\end{itemize}

\subsection{推荐的 Tag 命名规范}

\begin{verbatim}
# 论文版本
paper-<venue>-<version>
例如：paper-icml2026-v1, paper-icml2026-revision

# 实验组
exp-<experiment-name>
例如：exp-ablation-study, exp-baseline-comparison

# 主要结果
result-<result-name>
例如：result-table3-main, result-fig2-comparison

# 里程碑
milestone-<description>
例如：milestone-first-sota, milestone-reproducible-baseline
\end{verbatim}

\subsection{Tag 使用实践}

\paragraph{为论文的每个重要实验打 tag：}

\begin{verbatim}
# 跑完主实验后立即打 tag
git tag -a result-main-experiment -m \
  "Main results reported in Table 2, config: configs/main.yaml"

# 记录关键信息在 tag message 里
git tag -a result-ablation-study -m \
  "Ablation study results (Table 3)
   Run IDs: 2026-02-01_1030_ablation_*
   Config: configs/ablation_*.yaml
   Key finding: attention mechanism contributes 5% improvement"
\end{verbatim}

\paragraph{复现时直接切换到 tag：}

\begin{verbatim}
# 查看所有实验相关的 tag
git tag -l "result-*"

# 切换到特定实验版本
git checkout result-main-experiment

# 复现实验
make reproduce CONFIG=configs/main.yaml
\end{verbatim}

\section{实验产物不进 Git：用 .gitignore 保持仓库干净}

\textbf{核心原则：}Git 管理源代码和配置，不管理实验产物。

\subsection{什么不应该提交到 Git}

\begin{itemize}
  \item \textbf{模型权重：}通常很大（几百 MB 到几 GB），用专门的模型管理工具（如 DVC、Git LFS、或云存储）。
  \item \textbf{训练日志：}outputs/ 下的所有运行产物，按 run\_id 组织后归档或清理。
  \item \textbf{中间数据：}缓存的特征、预处理结果等，应该可重新生成。
  \item \textbf{数据集：}原始数据通常由外部管理，只在 data/ 里放小样本或数据指针（清单、下载脚本）。
  \item \textbf{虚拟环境：}venv/、.conda/ 等目录，用 requirements.txt 或 environment.yaml 代替。
\end{itemize}

\subsection{推荐的 .gitignore 模板}

\begin{verbatim}
# Python
__pycache__/
*.py[cod]
*$py.class
*.so
.Python

# 虚拟环境
venv/
env/
.conda/

# 实验产物
outputs/
runs/
checkpoints/
*.pt
*.pth
*.ckpt
*.h5

# 数据（除非是小样本）
data/raw/
data/processed/
*.csv
*.parquet

# 日志
*.log
logs/
wandb/

# 临时文件
.DS_Store
*.swp
*.swo
*~

# IDE
.vscode/
.idea/
*.iml

# 例外：保留小样本数据和配置
!data/samples/
!configs/
\end{verbatim}

\section{常见问题与解决方案}

\subsection{Q1：代码已经改了很多，怎么补救？}

如果你的仓库历史已经很混乱，不要试图"重写历史"（除非你非常熟悉 Git rebase）。推荐的做法：

\begin{enumerate}
  \item \textbf{设定基准点：}在当前状态打一个 tag：\texttt{git tag baseline-before-cleanup}
  \item \textbf{从现在开始规范：}
  \begin{itemize}
    \item 每个新实验使用独立分支
    \item 每个 commit 保持原子性
    \item 重要结果立即打 tag
  \end{itemize}
  \item \textbf{历史问题逐步修复：}
  \begin{itemize}
    \item 识别出论文关键实验对应的代码版本，补打 tag
    \item 在 README 或文档里记录"历史版本对应关系"
    \item 新实验都用规范流程，旧实验尽力追溯
  \end{itemize}
\end{enumerate}

\subsection{Q2：团队协作时如何统一分支策略？}

\begin{itemize}
  \item \textbf{写进 README：}把分支命名规范、tag 使用方法写进文档。
  \item \textbf{设定保护规则：}在 GitHub/GitLab 上设置 main 分支保护，禁止直接 push，必须通过 PR/MR。
  \item \textbf{Code Review：}合并到 main 前，检查是否满足 DoD（第5章），是否有完整的实验记录。
  \item \textbf{定期清理：}每周一次会议，集体清理无用的实验分支，归档重要的 tag。
\end{itemize}

\subsection{Q3：如何处理"dirty" 状态的实验？}

有时你临时修改了代码跑实验，还没来得及 commit，这就是"dirty"状态。

\textbf{记录策略：}
\begin{itemize}
  \item 在 run.json 里记录 \texttt{"git\_dirty": true}
  \item 同时记录 diff：\texttt{git diff > outputs/<run\_id>/changes.patch}
  \item 在 run.md 里注明临时改动的内容和原因
\end{itemize}

\textbf{事后补救：}
\begin{itemize}
  \item 如果结果有价值，立即把改动 commit 并打 tag
  \item 如果只是临时试验，在 run.md 里记录即可，不必 commit
\end{itemize}

\section{实战案例：从混乱到清晰的 Git 历史}

\subsection{重构前（反面教材）}

\begin{verbatim}
* a3f2d1c (HEAD -> main) update
* f8d9e0a fix
* 1b2c3d4 add new feature
* 9e8d7f6 initial commit
\end{verbatim}

无法从历史看出任何有用信息，所有论文实验找不到对应版本。

\subsection{重构后（正面示例）}

\begin{verbatim}
* d1e2f3g (tag: paper-icml2026-v1, main) Merge exp/final-ablation
|           Paper results ready for submission
|\
| * c4d5e6f (exp/final-ablation) Add ablation study for attention
| * b3c4d5e Configure ablation experiments
|/
* a1b2c3d (tag: result-main-experiment) Main experiment: achieve 95.2% accuracy
           Run ID: 2026-02-01_1030_main_run
           Config: configs/main_experiment.yaml
* 9a8b7c6 (tag: milestone-baseline) Establish reproducible baseline
           All baseline experiments validated
* 8f7e6d5 Fix data preprocessing bug in train/val split
* 7e6d5c4 Add comprehensive smoke test
* 6d5c4b3 Refactor data loading module
\end{verbatim}

清晰的历史，每个重要节点都有标记，可以随时回溯。

\section{10 分钟动作：为当前项目建立 Git 基线}

如果你现在只做一件事：为你的项目建立一个清晰的 Git 基线。

\begin{enumerate}
  \item \textbf{检查当前状态：}
  \begin{verbatim}
  git status
  git log --oneline -10
  \end{verbatim}

  \item \textbf{如果有未提交改动，决定处理方式：}
  \begin{itemize}
    \item 有价值的改动：整理后 commit，写清楚 message
    \item 临时试验：记录到 run.md，然后 \texttt{git stash}
    \item 无用改动：\texttt{git checkout .} 撤销
  \end{itemize}

  \item \textbf{为当前稳定版本打基线 tag：}
  \begin{verbatim}
  git tag -a baseline-$(date +%Y%m%d) -m \
    "Baseline before implementing git workflow"
  \end{verbatim}

  \item \textbf{设置规范的 .gitignore：}
  \begin{verbatim}
  # 使用前面提供的模板
  curl -o .gitignore <模板链接>
  # 或手动创建
  git add .gitignore
  git commit -m "Add comprehensive .gitignore for research project"
  \end{verbatim}

  \item \textbf{写下分支命名规范：}在 README.md 里添加一节"Git Workflow"，记录：
  \begin{itemize}
    \item main 分支用途
    \item exp/ 分支命名规范
    \item tag 使用方法
  \end{itemize}
\end{enumerate}

从现在开始，每个新实验都按规范使用分支和 tag，让 Git 真正成为你的"历史证明工具"。
